% !TEX root = ../termpaper.tex
% @author Leonhard Lüdeke
%
\newpage
\section{Lösungsansätze}
Wie oben beschrieben ist es, aufgrund des großen Suchraumes, unpraktikabel alle Spielstände abzubilden. Der Unterschied, der einzelnen Lösungsansätze, liegt alleine darin wie dieser Zustandsraum durchsucht wird und wie der Suchbaum aufgebaut wird. Alle Verfahren haben dabei gemein, dass sie den Zustandsraum mit Knoten (Stellungen im Spiel) und Kanten (ein Zug) darstellen. Nachfolgend sind Knoten und Kanten als diese Zustände zu verstehen. Je nachdem welcher Algorithmus beschrieben wird, können diese Zustände zusätzliche Eigenschaften bekommen. 

\subsection{Minimax-Algorithmus}
Ein fundamentales Konzept zur Berechnung von optimalen Strategien in Nullsummenspielen ist der Minimax-Algorithmus \cite{Solr-627596169}. Dieser Algorithmus basiert auf einem Spielbaum, in dem jeder Knoten ein MAX- oder MIN-Knoten ist. Diese geben, neben dem Spielzustand, an welcher Person (Maxi- oder Minimierer/in) am Zug ist und besitzen einen numerischen Wert. Dieser numerische Wert gibt an welcher Spielzustand besser für die jeweilige Person ist. Der sogenannte Maximierer/in (meist 1. Spieler/in) wählt immer den Spielzustand mit dem höheren wert und der Minimierer/in (meist 2. Spieler/in), den kleineren Wert. Dies entspricht einer Optimalen Spielweise. Werte gegen Null deuten auf eine ausgeglichenden Spielsituation hin, negative Werte gereichen zum Vorteil für Minimierer/in und positive Werte für Maximierer/in.

Die Werte haben ihren Ursprung in den Blattknoten (Endzustände) des Baumes. Diese besitzen eine Bewertungsfunktion, die aussagt wie gut oder schlecht dieser Endzustand für die Maximierende Person ist. Somit setzen sich die Werte für die Inneren Knoten aus Ihren Kindknoten zusammen, also durch das Propagieren der Werte Ihrer Kindknoten. Dabei gilt für die Maximierenden Knoten $Wert =$ $Maximum$ $aller$ $Kindwerte$ und für die Minimierenden Knoten $Wert =$ $Minimun$ $aller$ $Kindwerte$. Die Vorgehensweise mittels negativ Bewertungen wird die Negamax-Variante genannt. \newpage

Klassisch wird dieser Algorithmus mittels eine Tiefensuche implementiert und besitzt ohne Optimierung eine exponentielle Laufzeit von $O(b^d)$ wobei $b=moegliche \ Spalten \ pro \ Zug$ und $d=Zuege \ im \ Vorraus$ darstellen. Aufgrund dieser Großen Laufzeit ist eine Verbesserung notwendig, solch eine Optimierung bietet die Alpha-Beta Pruning Technik. Dabei werden die Äste des Spielbaums abgeschnitten die nicht für eine Optimale Lösung benötigt werden \cite{KNUTH1975293}. 

\subsection{Monte-Carlo-Treesearch}
Ein alternatives Konzept zur Berechnung von annähernd optimalen Strategien in Nullsummenspielen ist die Monte-Carlo-Baumsuche (MCTS). Wie der Minimax-Algorithmus basiert auch MCTS auf einem Spielbaum, in dem jeder Knoten einen Spielzustand repräsentiert und einen numerischen Wert besitzt. Im Gegensatz zum Minimax-Algorithmus, der durch statische Bewertungsfunktionen an den Blattknoten operiert, bestimmt MCTS diese Werte durch wiederholte stochastische Simulationen (Rollouts) bis zu den Endzuständen des Spielbaums. Durch die Propagierung der Simulationsergebnisse von den Blattknoten zu inneren Knoten konvergiert die Schätzung der Knotenwerte mit hinreichend vielen Iterationen gegen gute Näherungen der optimalen Spielqualität \cite{Kocsis2006}.

Die MCTS-Methode durchläuft zur Wertberechnung vier konzeptionelle Phasen: In der Selection-Phase werden bestehende Knoten gemäß der UCB-Heuristik ausgewählt, die Exploration und Exploitation balanciert. Die Expansion-Phase fügt neue Knoten hinzu, während die Simulation-Phase Zufallsspielzüge bis zum Endzustand durchführt. In der abschließenden Backpropagation-Phase werden die Ergebnisse dieser Simulationen entlang des durchlaufenen Pfades rückwärts propagiert und aktualisieren die Knotenwerte sowie Besuchszahlen.

Diese Methode wird iterativ implementiert und benötigt, im Gegensatz zum klassischen Minimax mit Tiefensuche, keine explizite Bewertungsfunktion. Während Minimax ohne Optimierung exponentielle Laufzeit von $O(b^d)$ aufweist, wobei $b = moegliche \ Spalten \ pro \ Zug$ und $d = Zuege im Voraus$, funktioniert MCTS praktikabler bei Spielen mit moderater Komplexität wie 4 Gewinnt. Mit nur 32 MB Speicher und etwa 500.000 Simulationen erreicht MCTS zuverlässige Spielstärke. Damit bietet MCTS eine effiziente Alternative zu Alpha-Beta Pruning, insbesondere wenn keine domänenspezifische Evaluierungsfunktion vorhanden ist.

